\chapter{Prefix sums, difference arrays, and two pointers}\label{ch:prefix}
\index{prefix sums}\index{difference array}\index{two pointers}

Chapters 1--4 built answers by searching, sorting, and arguing. This
chapter is quieter: three bookkeeping patterns that turn interval work
from $O(n)$ per query into $O(1)$ per query --- or a full scan into a
single sweep. The prefix sum precomputes every cumulative total once so
each later range-sum is one subtraction. The difference array is the
same idea run backwards: one range update is two point writes, and a
final prefix pass restores every cell. Two pointers replace a doubly
nested scan with one march when the window's contents are monotone.
After this chapter you will be able to answer a hundred thousand range
queries in linear time total, batch range updates for free, and
recognise exactly when the telescoping window is licensed.

% from-zero:begin
\begin{rulebox}[Before you start]
Chapters~1 and~4. \emph{Recap:} an array query ``sum of $a_l\dots a_r$'' naively costs $O(r-l+1)$ per query; with $q$ queries on $n$ elements that is $O(nq)$, which fails when both are $2\cdot10^5$. This chapter removes the repeated work.
\end{rulebox}
\begin{wins}[After this chapter you can]
\begin{itemize}
\item answer any number of range-sum queries in $O(1)$ each after one $O(n)$ pass;
\item apply many range updates at once and read the final array with a difference array;
\item extend both ideas to a grid with inclusion--exclusion;
\item solve ``longest/shortest subarray with property $P$'' with two pointers instead of two nested loops.
\end{itemize}
\end{wins}
\begin{principle}[How to think about this]
If you keep recomputing the same thing, store it. A prefix sum is the simplest case: spend one pass remembering running totals, and every later question becomes a subtraction. Whenever a problem repeats ``over a range'' or ``over a window,'' ask: \emph{what small summary could I build once so that each query needs only a glance?}
\end{principle}
% from-zero:end

\section{Prefix sums: build once, query forever}\label{sec:prefix}
\index{prefix sums!1D}

Given an array $a_1, \dots, a_n$, the \emph{prefix table} $P_i = a_1 +
\dots + a_i$ (with $P_0 = 0$) stores every partial sum in one pass of
adds. Any later range query is one lookup and one subtraction:
$\sum_{i=l}^{r} a_i = P_r - P_{l-1}$. The whole bargain is $O(n)$ build,
$O(1)$ query, forever --- the oldest space-for-time trade in the book,
and the first thing to try whenever a statement repeats ``the sum over
a range'' more than once.

\begin{lstlisting}[language=C++,caption={\texttt{prefix.h}: 1D and 2D
prefix tables.},label={lst:prefix}]
\end{lstlisting}
\lstinputlisting{code/prefix.h}

\begin{example}[CSES \emph{Static Range Sum Queries}]\label{ex:rangesum}
$n$ values, $q$ queries ``sum over $[a,b]$'', $n,q \le 2\cdot10^5$.
\par\smallskip
\textbf{Strategy.} $q$ independent sums, $O(n)$ each, is $4\cdot10^{10}$
adds --- dead. The constraint isn't the input size, it's the repetition:
precompute $P$, then every query is one subtraction.
\par\smallskip
\begin{solution}
\lstinputlisting{code/range_sum.cpp}
Real output (CSES sample $\{3,2,4,5,1,1,5,3\}$, queries
$[2,4],[5,6],[1,8],[3,3]$):
\begin{lstlisting}[numbers=none]
11
2
24
4
\end{lstlisting}
\end{solution}
\end{example}

\begin{figure}[ht]\centering
% fig19-prefixsum.tex -- Figure 5.1: prefix cells -> range sum as subtraction.
% fstyle.tex -- shared TikZ styles for the figures of this book (palette from lifeos.sty).
\tikzset{%
  cell/.style={draw=ink, minimum width=7mm, minimum height=6.5mm, inner sep=0pt, font=\footnotesize, fill=white},
  cellh/.style={cell, fill=accentsoft, draw=accent},
  cella/.style={cell, fill=amberbg, draw=amber},
  cellg/.style={cell, fill=goodbg, draw=good},
  cellx/.style={cell, fill=graybg, draw=slate, text=slate},
  vtx/.style={circle, draw=ink, minimum size=6.5mm, inner sep=0pt, font=\footnotesize, fill=white},
  vtxh/.style={vtx, draw=accent, fill=accentsoft},
  vtxa/.style={vtx, draw=amber, fill=amberbg},
  vtxg/.style={vtx, draw=good, fill=goodbg},
  vtxd/.style={vtx, draw=slate, fill=graybg, text=slate},
  lab/.style={font=\footnotesize, text=slate},
  labA/.style={font=\footnotesize, text=accent},
  labM/.style={font=\footnotesize, text=amber},
  idx/.style={font=\tiny, text=slate},
  arr/.style={-{stealth}, thick, ink},
  arrA/.style={-{stealth}, thick, accent},
  arrM/.style={-{stealth}, thick, amber},
  arrG/.style={-{stealth}, thick, good},
  edge/.style={thick, ink},
  edgeA/.style={thick, accent},
  edgeM/.style={thick, amber},
}%
% \cellrow{prefix}{x0}{y}{style}{v0,v1,...}: a row of cells named prefix0, prefix1, ... (pitch 0.75cm)
\providecommand{\cellrow}[5]{\foreach \v [count=\i from 0] in {#5} {\node[#4] (#1\i) at (#2+\i*0.75,#3) {\v};}}%
% \cellidx{prefix}{n}{start}: index labels under cells prefix0 .. prefix(n-1)
\providecommand{\cellidx}[3]{\foreach \i in {0,...,\numexpr#2-1\relax} {\node[idx, below=0pt of #1\i] {\the\numexpr\i+#3\relax};}}%
% \fitfig{<tikzpicture>}: shrink the picture to \linewidth only when it would overflow.
\ifdefined\fitbox\else\newsavebox{\fitbox}\fi%
\providecommand{\fitfig}[1]{\sbox{\fitbox}{#1}\ifdim\wd\fitbox>\linewidth\resizebox{\linewidth}{!}{\usebox{\fitbox}}\else\usebox{\fitbox}\fi}%

\begin{tikzpicture}[font=\footnotesize]
\node[lab, anchor=east] at (-0.4,0) {$a$};
\cellrow{a}{0.4}{0}{cell}{3,2,4,5,1,1}
\foreach \i [count=\k from 1] in {0,...,5} \node[idx, below=0pt of a\i] {$a_\k$};
\node[lab, anchor=east] at (-0.4,1.4) {$P$};
\foreach \i/\v in {0/0,1/3,2/5,3/9,4/14,5/15,6/16}
  \node[cellh, minimum width=7mm] (p\i) at (\i*0.75+0.025,1.4) {\v};
\foreach \i in {0,...,6} \node[idx, above=0pt of p\i] {$P_\i$};
\foreach \i in {1,2,3,4} \node[cell, fill=amberbg, draw=amber, minimum width=7mm] at (a\i) {};
\foreach \i/\v in {1/2,2/4,3/5,4/1} \node[font=\footnotesize] at (a\i) {\v};
\draw[decorate,decoration={brace,mirror,raise=11pt},amber,thick] (a1.south west) -- (a4.south east) node[midway,below=15pt,labM]{query $[2,5]$};
\node[cella, minimum width=7mm] at (p1) {3};
\node[cella, minimum width=7mm] at (p5) {15};
\node[labM, align=left, anchor=west] at (5.5,0.7) {sum $=P_5-P_1$\\$\phantom{\text{sum }}=15-3=12$\\[2pt]one subtraction,\\not four reads};
\end{tikzpicture}

\caption{Prefix cells versus a range query. Notice the query reads
\emph{two} of the blue cells, none of the data cells: $P_r - P_{l-1}$
subtracts everything before the range at once.}
\label{fig:prefixsum}
\end{figure}

\begin{exercise}[easy]\label{exr:prefixcalc}
Given $a = \{4, 1, 7, 2, 6\}$, build $P$ by hand and evaluate
$\sum_{3}^{5}$ and $\sum_{1}^{4}$ with single subtractions.
\end{exercise}

\section{Difference arrays: the inverse direction}\label{sec:diff}
\index{difference array!1D}

The prefix sum reads a finished array per range. Its inverse writes:
to add $v$ to every $a_i$ for $i \in [l, r]$, do \emph{not} touch the
range --- write $d_l \mathrel{+}= v$ and $d_{r+1} \mathrel{-}= v$ into
a difference array $d$. Two point writes per update, and a single
prefix pass of $d$ at the end restores every cell: position $i$'s
restored value is exactly the sum of all updates that started at or
before $i$ minus those that already closed. All updates batched, $O(1)$
each.

\begin{lstlisting}[language=C++,caption={\texttt{diff.h}: range updates
as two point writes, 1D and 2D.},label={lst:diff}]
\end{lstlisting}
\lstinputlisting{code/diff.h}

\begin{lstlisting}[numbers=none,caption={\texttt{diff\_demo.cpp} real
output --- three updates, one restore.},label={lst:diffout}]
updates: [2,4]+3, [1,3]+5, [6,7]+2
restored: 0 5 8 8 3 0 2 2
\end{lstlisting}

Check the restore by hand: cell 1 sees only $+5$; cell 3 sees
$+3{+}5$; cell 5 sits outside every update. Figure~\ref{fig:diffpipe}
shows the pipeline.

\begin{figure}[ht]\centering
\begin{tikzpicture}[font=\footnotesize,x=0.6cm]
% F5.4 diff pipeline: [2,4]+3 -> writes -> restore.
\node[anchor=east] at (-0.3,0) {writes:};
\foreach \i/\v in {0/0,1/0,2/+3,3/0,4/0,5/$-3$,6/0,7/0}
  \node[draw=slate,minimum width=0.6cm,minimum height=0.45cm,font=\tiny]
    at (\i*0.68,0) {\v};
\node[font=\tiny] at (1.36,0.42) {$d[2]$};
\node[font=\tiny] at (3.4,0.42) {$d[5]$};
\node[anchor=east] at (-0.3,-0.9) {prefix:};
\foreach \i/\v in {0/0,1/0,2/3,3/3,4/3,5/0,6/0,7/0}
  \node[draw=accent,fill=accentsoft,minimum width=0.6cm,
        minimum height=0.45cm,font=\tiny] at (\i*0.68,-0.9) {\v};
\draw[-{stealth},thick,ink] (2.4,-0.15) -- (2.4,-0.6);
\end{tikzpicture}

\caption{A difference-array update becomes two point writes; the
prefix pass restores cells. Notice what the $-v$ write buys: the
increase \emph{ends}, which is what makes the restored plateau flat.}
\label{fig:diffpipe}
\end{figure}

\section{Two dimensions: rectangles by inclusion--exclusion}\label{sec:2d}
\index{prefix sums!2D}\index{inclusion--exclusion!rectangles}

The same subtraction generalises to grids. $P_{r,c}$ = sum of the
$r \times c$ top-left corner. A rectangle $(r_1, c_1)$--$(r_2, c_2)$
then costs four lookups with the PIE sign pattern:
$P_{r_2,c_2} - P_{r_1-1,c_2} - P_{r_2,c_1-1} + P_{r_1-1,c_1-1}$ ---
the two over-subtracted corner strips come back once, exactly the
inclusion--exclusion of Figure~\ref{fig:ie2d}. And the difference
array generalises identically: a rectangle $+v$ update is four corner
writes, restored by a 2D prefix pass.

\begin{example}[CSES \emph{Forest Queries}]\label{ex:forest}
Count trees in axis-aligned rectangles of an $n \times n$ grid,
$q \le 2\cdot10^5$ queries.
\par\smallskip
\textbf{Strategy.} ``How many inside a rectangle'' repeated $q$ times
is the 2D prefix pattern verbatim: trees are just values in
$\{0,1\}$, so build $P$ once and answer each rectangle in four reads.
\par\smallskip
\begin{solution}
\lstinputlisting{code/forest_queries.cpp}
Real output (CSES sample):
\begin{lstlisting}[numbers=none]
3
1
2
\end{lstlisting}
\end{solution}
\end{example}

\begin{figure}[ht]\centering
% fig20-ie2d.tex -- Figure 5.2: 2D inclusion--exclusion on a 4-row x 5-column grid; target = rows 2..4, cols 2..4.
% P_{r,c} = sum of all cells with row <= r and column <= c (rows counted downward from the top).
% fstyle.tex -- shared TikZ styles for the figures of this book (palette from lifeos.sty).
\tikzset{%
  cell/.style={draw=ink, minimum width=7mm, minimum height=6.5mm, inner sep=0pt, font=\footnotesize, fill=white},
  cellh/.style={cell, fill=accentsoft, draw=accent},
  cella/.style={cell, fill=amberbg, draw=amber},
  cellg/.style={cell, fill=goodbg, draw=good},
  cellx/.style={cell, fill=graybg, draw=slate, text=slate},
  vtx/.style={circle, draw=ink, minimum size=6.5mm, inner sep=0pt, font=\footnotesize, fill=white},
  vtxh/.style={vtx, draw=accent, fill=accentsoft},
  vtxa/.style={vtx, draw=amber, fill=amberbg},
  vtxg/.style={vtx, draw=good, fill=goodbg},
  vtxd/.style={vtx, draw=slate, fill=graybg, text=slate},
  lab/.style={font=\footnotesize, text=slate},
  labA/.style={font=\footnotesize, text=accent},
  labM/.style={font=\footnotesize, text=amber},
  idx/.style={font=\tiny, text=slate},
  arr/.style={-{stealth}, thick, ink},
  arrA/.style={-{stealth}, thick, accent},
  arrM/.style={-{stealth}, thick, amber},
  arrG/.style={-{stealth}, thick, good},
  edge/.style={thick, ink},
  edgeA/.style={thick, accent},
  edgeM/.style={thick, amber},
}%
% \cellrow{prefix}{x0}{y}{style}{v0,v1,...}: a row of cells named prefix0, prefix1, ... (pitch 0.75cm)
\providecommand{\cellrow}[5]{\foreach \v [count=\i from 0] in {#5} {\node[#4] (#1\i) at (#2+\i*0.75,#3) {\v};}}%
% \cellidx{prefix}{n}{start}: index labels under cells prefix0 .. prefix(n-1)
\providecommand{\cellidx}[3]{\foreach \i in {0,...,\numexpr#2-1\relax} {\node[idx, below=0pt of #1\i] {\the\numexpr\i+#3\relax};}}%
% \fitfig{<tikzpicture>}: shrink the picture to \linewidth only when it would overflow.
\ifdefined\fitbox\else\newsavebox{\fitbox}\fi%
\providecommand{\fitfig}[1]{\sbox{\fitbox}{#1}\ifdim\wd\fitbox>\linewidth\resizebox{\linewidth}{!}{\usebox{\fitbox}}\else\usebox{\fitbox}\fi}%

\begin{tikzpicture}[font=\footnotesize, x=0.42cm, y=0.42cm]
% \mini{xshift}{shaded x0}{x1}{y0}{y1}{fill}
\newcommand{\mini}[6]{\begin{scope}[shift={(#1,0)}]
  \fill[#6] (#2,#4) rectangle (#3,#5);
  \foreach \x in {0,...,5} \draw[slate] (\x,0) -- (\x,4);
  \foreach \y in {0,...,4} \draw[slate] (0,\y) -- (5,\y);
  \draw[accent, thick, dashed] (1,0) rectangle (4,3);
  \end{scope}}
\mini{0}{0}{4}{0}{4}{accent!18}      \node at (6.1,2) {\Large $-$};
\mini{7}{0}{4}{3}{4}{amber!25}          \node at (13.1,2) {\Large $-$};
\mini{14}{0}{1}{0}{4}{amber!25}         \node at (20.1,2) {\Large $+$};
\mini{21}{0}{1}{3}{4}{good!30}          \node at (27.1,2) {\Large $=$};
\mini{28}{1}{4}{0}{3}{accent!18}
\node[font=\scriptsize] at (2.5,-1.1) {$P_{r_2,c_2}$};
\node[font=\scriptsize] at (9.5,-1.1) {$P_{r_1-1,\,c_2}$};
\node[font=\scriptsize] at (16.5,-1.1) {$P_{r_2,\,c_1-1}$};
\node[font=\scriptsize] at (23.5,-1.1) {$P_{r_1-1,\,c_1-1}$};
\node[font=\scriptsize, text=accent] at (30.5,-1.1) {target (dashed)};
\node[lab, align=center] at (15,-3.3) {the top-left strip is subtracted twice, so it is added back once --- the signs alternate $+\,-\,-\,+$};
\end{tikzpicture}

\caption{Rectangle sums by corner algebra. Notice the signs
$+,-,-,+$: the count over the target rectangle is the big corner
minus the two flanking strips, plus the little corner both strips
stole.}
\label{fig:ie2d}
\end{figure}

\begin{example}[USACO-S \emph{Painting the Barn}]\label{ex:barn}
\index{Painting the Barn}
$K$ coats visible: given a list of rectangles painted one coat at a
time, count the unit squares covered by \emph{exactly} $K$
rectangles (grid $\le 10^3$).
\par\smallskip
\textbf{Strategy.} The statement hands you an \emph{update list}, not
an array --- that is the difference array's native habitat. Every
rectangle is one 2D-range $+1$ update ($4$ corner writes); after all
updates, restore the grid once and count cells $= K$.
\par\smallskip
\begin{solution}
\lstinputlisting{code/painting_barn.cpp}
Real output ($N{=}5$, $K{=}3$, rectangles
$(1,1){-}(4,4)$, $(2,2){-}(4,4)$, $(3,2){-}(4,4)$ --- cells $(3,2)$
and $(3,3)$ carry all three coats):
\begin{lstlisting}[numbers=none]
2
\end{lstlisting}
\end{solution}
\end{example}

\section{Two pointers: the telescoping window}\label{sec:twoptr}
\index{two pointers}

Different machine, same moral: if a window only needs to slide, two
indices can do what a double loop did. For ``count subarrays summing
to $x$'' over \emph{positive} values: extend $r$ while the window is
short of $x$, contract $l$ while it overshoots. Each index moves only
rightward, so the whole march is $O(n)$ --- and each subarray counted
is genuinely the only window at that $l$ with the right sum, because
growing or shrinking a positive window monotonically changes the sum.

\begin{example}[CSES \emph{Subarray Sums I}]\label{ex:sums1}
Count subarrays summing to $x$, all $a_i$ positive.
\par\smallskip
\textbf{Strategy.} Positivity is the licence: window sums move
\emph{monotonically} as $l$ or $r$ advances, so each $l$ has at most
one valid $r$ to find. March, don't scan.
\par\smallskip
\begin{solution}
\lstinputlisting{code/twoptr_sums.cpp}
Real output (CSES sample $\{2,4,1,2,7\}$, $x{=}7$ --- windows
$\{4,1,2\}$, $\{7\}$, $\{2,7\}$):
\begin{lstlisting}[numbers=none]
3
\end{lstlisting}
\end{solution}
\end{example}

\begin{figure}[ht]\centering
% fig21-twoptr.tex -- Figure 5.3: two-pointer march on {2,4,1,2,7}, x=7 target sum 7; four frames of the march.
% fstyle.tex -- shared TikZ styles for the figures of this book (palette from lifeos.sty).
\tikzset{%
  cell/.style={draw=ink, minimum width=7mm, minimum height=6.5mm, inner sep=0pt, font=\footnotesize, fill=white},
  cellh/.style={cell, fill=accentsoft, draw=accent},
  cella/.style={cell, fill=amberbg, draw=amber},
  cellg/.style={cell, fill=goodbg, draw=good},
  cellx/.style={cell, fill=graybg, draw=slate, text=slate},
  vtx/.style={circle, draw=ink, minimum size=6.5mm, inner sep=0pt, font=\footnotesize, fill=white},
  vtxh/.style={vtx, draw=accent, fill=accentsoft},
  vtxa/.style={vtx, draw=amber, fill=amberbg},
  vtxg/.style={vtx, draw=good, fill=goodbg},
  vtxd/.style={vtx, draw=slate, fill=graybg, text=slate},
  lab/.style={font=\footnotesize, text=slate},
  labA/.style={font=\footnotesize, text=accent},
  labM/.style={font=\footnotesize, text=amber},
  idx/.style={font=\tiny, text=slate},
  arr/.style={-{stealth}, thick, ink},
  arrA/.style={-{stealth}, thick, accent},
  arrM/.style={-{stealth}, thick, amber},
  arrG/.style={-{stealth}, thick, good},
  edge/.style={thick, ink},
  edgeA/.style={thick, accent},
  edgeM/.style={thick, amber},
}%
% \cellrow{prefix}{x0}{y}{style}{v0,v1,...}: a row of cells named prefix0, prefix1, ... (pitch 0.75cm)
\providecommand{\cellrow}[5]{\foreach \v [count=\i from 0] in {#5} {\node[#4] (#1\i) at (#2+\i*0.75,#3) {\v};}}%
% \cellidx{prefix}{n}{start}: index labels under cells prefix0 .. prefix(n-1)
\providecommand{\cellidx}[3]{\foreach \i in {0,...,\numexpr#2-1\relax} {\node[idx, below=0pt of #1\i] {\the\numexpr\i+#3\relax};}}%
% \fitfig{<tikzpicture>}: shrink the picture to \linewidth only when it would overflow.
\ifdefined\fitbox\else\newsavebox{\fitbox}\fi%
\providecommand{\fitfig}[1]{\sbox{\fitbox}{#1}\ifdim\wd\fitbox>\linewidth\resizebox{\linewidth}{!}{\usebox{\fitbox}}\else\usebox{\fitbox}\fi}%

\fitfig{\begin{tikzpicture}[font=\footnotesize]
% \frame{x}{l}{r}{caption}
\newcommand{\fr}[4]{  \foreach \v [count=\i from 0] in {2,4,1,2,7}
    \node[\ifnum\i<#2 cellx\else\ifnum\i>#3 cellx\else cellh\fi\fi, minimum width=6.5mm] (w\i) at (#1+\i*0.65,0) {\v};
  \node[labA, below=1pt of w#2] {$l$};
  \node[labA, above=1pt of w#3] {$r$};
  \node[lab, below=14pt of w2, text width=3cm, align=center] {#4};
}
\fr{0}{0}{1}{sum $=6$: grow $r$}
\fr{3.9}{0}{2}{sum $=7$: count, then grow}
\fr{7.8}{1}{3}{sum $=7$: count, shrink}
\fr{11.7}{1}{4}{sum $=14$: shrink $l$}
\end{tikzpicture}}

\caption{The two-pointer march on the CSES sample. Notice each pointer
travels in one direction only --- $l$ never revisits an index, which is
where the $n^2$ went.}
\label{fig:twoptr}
\end{figure}

\begin{exercise}[easy]\label{exr:positivity}
The two-pointer code counts \emph{wrongly} once negatives are allowed.
Explain which step of the march breaks and why (find a counterexample
with a single $-1$).
\end{exercise}

\section{Sliding windows with bookkeeping}\label{sec:window}
\index{sliding window}\index{prefix sums!hash counting}

When values can be negative, the window stops telescoping --- but the
prefix table still answers, this time by \emph{counting pairs}:
subarray $[i{+}1,j]$ has sum $x$ iff $P_j - P_i = x$. Walk $j$ once,
counting how many earlier prefixes equal $P_j - x$ in a \texttt{map}.
The same trick with remainders counts subarrays divisible by $n$:
equal prefix residues. Both are $O(n)$ or $O(n \log n)$ where the
telescoping window was forbidden.

\begin{example}[CSES \emph{Subarray Sums II} + \emph{Subarray
Divisibility}]\label{ex:sums2}\index{Subarray Sums II}
Same count with negatives allowed; and subarrays whose sum is a
multiple of $n$.
\par\smallskip
\textbf{Strategy.} Negatives kill monotonicity, so switch the
predicate's frame: the window condition $P_j - P_i = x$ is a
\emph{pair} condition between prefix values --- count prefix
occurrences, not windows. For divisibility, the condition becomes
equal residues mod $n$.
\par\smallskip
\begin{solution}
\lstinputlisting{code/prefix_map.cpp}
Real output (CSES samples $\{2,-1,3,5,-2\}$, $x{=}7$;
$\{3,1,2,7,4\}$ divisibility):
\begin{lstlisting}[numbers=none]
sum=7: 2
divisible by n=5: 1
\end{lstlisting}
\end{solution}
\end{example}

\begin{figure}[ht]\centering
\begin{tikzpicture}[font=\footnotesize]
% F5.5 residue buckets for {3,1,2,7,4} mod 5.
\foreach \i/\lbl in {0/r{=}0,1/r{=}1,2/r{=}2,3/r{=}3,4/r{=}4} {
  \node[draw=slate,minimum width=1.1cm,minimum height=0.5cm,font=\tiny]
    (b\i) at (\i*1.3,0) {\lbl};
}
\node[font=\tiny,above] at (b0) {$P_0$};
\node[font=\tiny,above] at (b1) {$P_3$};
\node[font=\tiny,above] at (b2) {$P_5$};
\node[font=\tiny,above] at (b3) {$P_1, P_4$};
\node[font=\tiny,above] at (b4) {$P_2$};
\node[align=left,text width=4.6cm] at (8.6,0.3)
  {prefixes $P_i$ by residue:\\
   bucket 3 has a \emph{pair}\\
   $P_1{=}P_4{=}3$ $\Rightarrow$\\
   $a_{2..4}$ sums to $0 \bmod 5$};
\end{tikzpicture}

\caption{Prefix-remainder counting on $\{3,1,2,7,4\}$ mod $5$.
Notice the answer comes from a \emph{repeated residue}, not from any
single value: two equal prefix sums bracket a divisible subarray.}
\label{fig:buckets}
\end{figure}

\begin{lstlisting}[numbers=none,caption={\texttt{window\_demo.cpp}:
sliding distinct-count and deque minimum preview.},label={lst:window}]
window[0..2] distinct=2 min=1
window[1..3] distinct=3 min=1
window[2..4] distinct=3 min=1
window[3..5] distinct=2 min=2
\end{lstlisting}

The deque line is a preview: Chapter~6 formalises the monotonic queue
that answers ``window minimum'' in amortised $O(1)$.

\begin{exercise}[medium]\label{exr:remainder}
Show that if two prefix sums are congruent mod $m$, the subarray
between them is divisible by $m$. Is the converse true?
\end{exercise}

\section{Chapter summary}
\begin{table}[ht]\centering
{\small
\begin{tabularx}{\linewidth}{@{}lYYY@{}}
\toprule
\tablehead\thead{Asked} & \thead{Tool} & \thead{Cost} & \thead{Needs}\\
\midrule
sum over $[l,r]$, $q$ times & 1D prefix $P$ & $O(n)+O(1)$ & static array\\
rectangle sums & 2D prefix & 4 lookups & static grid\\
range $+v$ updates & difference array & 2 writes + 1 pass & batch offline\\
subarray sum $=x$, $a_i>0$ & two pointers & $O(n)$ & positivity!\\
subarray sum $=x$, any $a_i$ & prefix + \texttt{map} & $O(n\log n)$ & pair count\\
divisible-sum subarrays & prefix residues & $O(n)$ & mod $n$\\
window stats & deque/\texttt{map} slide & $O(n)$ & Ch~6 preview\\
\bottomrule
\end{tabularx}}
\caption{Every row is the same idea at a different orientation:
precompute cumulative structure once, answer forever.}
\label{tab:prefix}
\end{table}

\section{Common mistakes}
\index{common mistakes}
\begin{itemize}
\item \textbf{Two pointers on negative data.} The window's sum is not
monotone in $l$ once a negative value exists, so the march skips
valid answers --- Subarray Sums II exists precisely because the
telescope fails. Reach for prefix-pair counting instead.
\item \textbf{$P_0$ forgetting.} A prefix table without the sentinel
$P_0 = 0$ cannot answer a range that starts at index $1$ --- every
book keeps the zero cell.
\item \textbf{Off-by-one on the difference close.} $d_{r+1}$ ends the
update, not $d_r$: the subtract write belongs one step \emph{after}
the range, so the restore is a prefix of $d$, never of $a$.
\item \textbf{Rectangle corner confusion.} Half-open vs closed corners
disagree by one in both axes; fix the convention (ours: $P_{r,c}$ is
the $r\times c$ corner, query subtracts at $r_1-1$, $c_1-1$) and test
on the $1\times1$ case.
\item \textbf{Residues that go negative.} $P \bmod m$ in C++ keeps the
sign of $P$; always normalise $((P \% m) + m) \% m$ before bucketing.
\end{itemize}

\section{Additional exercises}
\begin{exercise}[easy]
USACO-S \emph{Breed Counting} --- three prefix tables, one per breed.
\end{exercise}
\begin{exercise}[easy]
USACO-S \emph{Hoof, Paper, Scissors} --- prefix counts per gesture.
\end{exercise}
\begin{exercise}[easy]
CSES \emph{Maximum Subarray Sum} --- prefix min so far, one pass.
\end{exercise}
\begin{exercise}[easy]
CF \emph{Karen and Coffee} --- difference array over the recipe range.
\end{exercise}
\begin{exercise}[medium]
CSES \emph{Subarray Sums I} --- implement the telescoping window.
\end{exercise}
\begin{exercise}[medium]
CSES \emph{Subarray Sums II} --- prefix + \texttt{map} counting.
\end{exercise}
\begin{exercise}[medium]
CSES \emph{Subarray Divisibility} --- residue buckets mod $n$.
\end{exercise}
\begin{exercise}[medium]
USACO-S \emph{Subsequences Summing to Sevens} --- first/last prefix
per residue.
\end{exercise}
\begin{exercise}[medium]
SPOJ \emph{Haybale Stacking} --- difference array then read off.
\end{exercise}
\begin{exercise}[medium]
CSES \emph{Distinct Values Subarrays} --- sliding window + last-seen.
\end{exercise}
\begin{exercise}[hard]
USACO-S \emph{Rectangular Pasture} --- 2D prefix over a sparse grid.
\end{exercise}
\begin{exercise}[hard]
USACO-G \emph{Haybale Feast} --- window + monotone queue.
\end{exercise}
\begin{exercise}[hard]
CF \emph{Queries for Number of Palindromes} --- a prefix table with a
cleverer cell content.
\end{exercise}

\chapternotes
The prefix/difference pair and the telescoping window follow
Sannemo's array-techniques sections~\cite{sannemo2025} and
Laaksonen's range-query chapter~\cite{laaksonen2024}; the USACO
Guide's Silver prefix-sum and two-pointer modules~\cite{usacoguide}
are the direct drill source, and Halim~\cite{halimcp4} covers the
same machinery with more worked IDs. The five archive problems are
official CSES/USACO tasks~\cite{cses}; the remainder-counting trick
appears throughout the cp-algorithms notes~\cite{cpalgorithms}.

\section*{Further reading}
\begin{itemize}
\item USACO Guide, ``Prefix Sums'' (intermediate).
\url{https://usaco.guide/silver/prefix-sums} --- 1-D and 2-D prefix sums
with more graded problems at exactly this chapter's level.
\item USACO Guide, ``Two Pointers'' (intermediate).
\url{https://usaco.guide/silver/two-pointers} --- the window idiom of
\S\ref{sec:twoptr} with a deeper problem bank.
\item Laaksonen, \emph{Competitive Programmer's Handbook}, ch.~8--9
(beginner). \url{https://cses.fi/book/book.pdf} --- prefix sums and the
sliding-window minimum-sum problems in one short read.
\item CSES 1646 \emph{Static Range Sum Queries} (practice).
\url{https://cses.fi/problemset/task/1646} --- the direct archive form
of \ref{ex:rangesum}.
\item CSES 1660 \emph{Subarray Sums I} and 1644 \emph{Maximum Subarray
Sum II} (practice). \url{https://cses.fi/problemset/task/1660} ---
live exercises for \S\ref{sec:twoptr} and \S\ref{sec:window}.
\end{itemize}

\begin{rulebox}[Watch instead]
\smallskip\noindent\textbf{Best first.}
``Prefix Sum Array and Range Sum Queries'' --- Profound Academy,
8 min --- \url{https://www.youtube.com/watch?v=xbYr9JOC2Lk}. The
accumulation trick of \S\ref{sec:prefix} animated cell by cell;
uses \texttt{sum[i]} where the book writes $p_i$ --- same array.
\smallskip\noindent
``Two Pointers in 7 minutes'' --- AlgoMasterIO, 7 min ---
\url{https://www.youtube.com/watch?v=QzZ7nmouLTI}. The
shrink-on-violation loop of \S\ref{sec:twoptr} on interview-style
problems; Python-flavoured, the pointer invariant is identical.
Access date 2026-10-10.
\end{rulebox}
